% ============================================================================
%  1. Introducción y objetivos
% ============================================================================
\section{Introducción y objetivos}
\label{sec:introduccion}

Un robot móvil diferencial se desplaza y orienta variando solo la velocidad y
el sentido de sus dos ruedas motrices: misma velocidad y sentido produce
traslación, sentidos opuestos producen rotación sobre su eje. Esta
arquitectura es la base de robots de servicio, AGV industriales y plataformas
de competencia, y exige controlar dos motores DC de forma independiente y
sincronizada.

\subsection*{Objetivo general}
Implementar el control de velocidad, sentido de giro y frenado de una
plataforma móvil diferencial con un PIC16F877A y un driver DRV8833, validado en
simulación y en hardware real.

\subsection*{Objetivos específicos}
\begin{enumerate}
  \item Configurar y sincronizar los módulos CCP1 (RC2) y CCP2 (RC1) en modo
        PWM con Timer2 como base de tiempo única.
  \item Integrar el DRV8833 para gobernar velocidad, sentido de giro y frenado
        dinámico de dos motores DC.
  \item Aplicar normas de montaje en protoboard: masa común, habilitación por
        nSLEEP y supresión de ruido inductivo ($dI/dt$).
  \item Desarrollar firmware modular en MPLAB XC8 para la navegación de la
        plataforma.
  \item \emph{(Extensión del grupo)} Agregar control remoto desde un celular
        mediante un módulo Bluetooth serie.
\end{enumerate}
